\chapter{System Requirement Specification}
\label{ch:srs}

\section{Functional Requirements}
\label{sec:functional}
The artificial intelligence-based pipeline optimization system should satisfy the following functional requirements. Section~\ref{sec:testing} maps each one to its implementation.
\begin{itemize}
\item Collection of historical CI/CD pipeline records, including build outcomes, commits, timing and logs, and live ingestion of GitHub Actions events.
\item Preprocessing and cleansing of pipeline data for effective analysis.
\item Feature extraction from historical records, including code churn, failure rates, commit timing and pipeline history, using only information available before a build runs.
\item Training and execution of machine learning algorithms for build failure prediction.
\item Anomaly detection for slow or abnormal pipeline jobs.
\item Categorization of similar failure logs for automated root cause analysis.
\item Optimization suggestions to improve the efficiency of pipelines.
\item Visualization of pipeline analytics, predictions and bottlenecks through a dashboard.
\end{itemize}

\section{Non-functional Requirements}
\label{sec:nonfunctional}
\begin{itemize}
\item \textbf{Scalability:} The system should cope efficiently with large volumes of CI/CD data from many projects and repositories. Development used 3.56 million Travis~CI job records.
\item \textbf{Reliability:} The platform should produce consistent predictions. The live history computed by the service must equal the history used in training.
\item \textbf{Performance:} The system should score a build quickly enough to be useful before it runs.
\item \textbf{Accuracy:} Models should rank risky builds well within each repository and be evaluated without information leakage.
\item \textbf{Explainability:} Each prediction should come with an explanation of the factors behind it.
\item \textbf{Security:} Webhook endpoints must verify GitHub's HMAC signature, and the public demonstration must be read-only.
\item \textbf{Usability:} The monitoring dashboard should be user-friendly and intuitive.
\item \textbf{Maintainability:} The architecture should make maintenance and enhancement easy; automated tests and continuous integration protect it.
\end{itemize}

\section{Hardware Requirements}
\label{sec:hardware}
\begin{table}[H]
\centering
\caption{Hardware requirements}
\label{tab:hw}
\begin{tabular}{|p{4cm}|p{9cm}|}
\hline
\textbf{Component} & \textbf{Specification}\\ \hline
Processor & Intel i5 / Ryzen 5 / Apple silicon or higher\\ \hline
RAM & Minimum 8 GB (16 GB recommended for training on the full Travis~CI dataset)\\ \hline
Storage & 256 GB SSD\\ \hline
GPU (Optional) & Not required: all models are tree ensembles trained on CPU\\ \hline
Operating System & Windows / Linux / macOS\\ \hline
Deployment & Cloud container service with 512 MB RAM (Render free tier is sufficient for the demonstration)\\ \hline
\end{tabular}
\end{table}

\section{Software Requirements}
\label{sec:software}
\begin{longtable}{|p{4.6cm}|p{8.4cm}|}
\caption{Software requirements}\label{tab:sw}\\
\hline
\textbf{Software Component} & \textbf{Purpose}\\ \hline
\endfirsthead
\hline
\textbf{Software Component} & \textbf{Purpose}\\ \hline
\endhead
Python 3.11 & Core programming language\\ \hline
Pandas, NumPy, PyArrow & Data preprocessing, analysis and storage\\ \hline
Scikit-learn & Random Forest, Decision Tree, SVM, Isolation Forest, TF-IDF and k-means\\ \hline
LightGBM, XGBoost & Gradient-boosted models, including the learning-to-rank model\\ \hline
SHAP & Explanations of individual predictions\\ \hline
FastAPI, Uvicorn & REST API and webhook service\\ \hline
SQLAlchemy, SQLite & Storage of pipelines, builds, predictions and job telemetry\\ \hline
ReportLab & Per-build PDF reports\\ \hline
React, Vite, Tailwind CSS & Web dashboard\\ \hline
Recharts & Animated charts on the dashboard\\ \hline
pytest & Automated tests (84 tests)\\ \hline
GitHub, GitHub Actions & Version control, continuous integration of this project, and a data source\\ \hline
Docker, Render & Containerized deployment of the public demonstration\\ \hline
Visual Studio Code & Development environment\\ \hline
\end{longtable}
